
\documentclass{article}
\usepackage{colortbl}
\usepackage{makecell}
\usepackage{multirow}
\usepackage{supertabular}

\begin{document}

\newcounter{utterance}

\centering \large Interaction Transcript for game `eqbench', experiment `eq\_bench', episode 133 with reflexion{-}Qwen3.5{-}27B.
\vspace{24pt}

{ \footnotesize  \setcounter{utterance}{1}
\setlength{\tabcolsep}{0pt}
\begin{supertabular}{c@{$\;$}|p{.15\linewidth}@{}p{.15\linewidth}p{.15\linewidth}p{.15\linewidth}p{.15\linewidth}p{.15\linewidth}}
    \# & \multicolumn{2}{c}{Player} && \multicolumn{2}{c}{Game Master} \\
    \hline

    \theutterance \stepcounter{utterance}  
    & & & \multicolumn{4}{p{0.6\linewidth}}{
        \cellcolor[rgb]{0.9,0.9,0.9}{
            \makecell[{{p{\linewidth}}}]{
                \texttt{\tiny{[P1$\langle$GM]}}
                \texttt{You have attempted this same task before and did not succeed. These are the plans you wrote after those attempts.} \\
\texttt{Trial \#1: New plan: My previous attempt failed because I assigned a high score (8) to "Defeated," which misinterprets Arthur's persistent, active efforts to apologize and reconcile as passive resignation. Arthur is not giving up; he is actively trying to bridge the gap despite Sam's rejection. To correct this, I will re{-}evaluate the dialogue by focusing on Arthur's repeated assertions of intent ("I'm trying to change," "I can try to make it up to you") rather than just the rejection he faces. I will significantly lower the score for "Defeated" to reflect his resilience and increase the weight on "Remorseful" as the primary driver of his actions, while keeping "Rebellious" and "Vengeful" at 0 as they contradict his conciliatory tone. The new strategy is to prioritize the emotion that motivates the character's continued engagement over the emotion that suggests withdrawal.} \\
\\ 
\texttt{Trial \#2: New plan: My previous attempt failed because I assigned a non{-}zero score (3) to "Defeated," which incorrectly interprets Arthur's persistent, active attempts to reconcile as passive resignation. Although Sam is rejecting him, Arthur continues to argue his case ("I'm trying to change," "I can try to make it up to you"), demonstrating resilience rather than defeat. To correct this, I will strictly evaluate the character's active engagement: since Arthur is still fighting for the relationship, "Defeated" must be scored as 0. I will maintain "Remorseful" as the highest score (9 or 10) to reflect his repeated apologies and acknowledgment of fault, while keeping "Rebellious" and "Vengeful" at 0 as they contradict his conciliatory tone. The new strategy is to assign a score of 0 to any emotion that implies withdrawal or aggression when the character is actively seeking resolution.} \\
\\ 
\texttt{Trial \#3: New plan: My previous attempt failed because I assigned a score of 0 to "Defeated," overlooking the emotional impact of Sam's final, absolute rejection ("Stop saying you're sorry. It doesn't change anything"). While Arthur is actively trying to reconcile (Remorseful), the futility of his efforts in the face of Sam's hardened stance suggests a secondary feeling of helplessness or defeat. To correct this, I will re{-}evaluate the ending not just for Arthur's intent, but for the outcome of his actions. I will assign a high score to "Remorseful" (10) as the primary emotion, but also assign a non{-}zero score to "Defeated" (e.g., 6 or 7) to reflect the crushing realization that his apologies are ineffective. "Rebellious" and "Vengeful" will remain 0. The new strategy is to balance the character's active intent with the passive emotional result of their failed attempts.} \\
\\ 
\\ 
\\ 
\texttt{Your task is to predict the likely emotional responses of a character in this dialogue:} \\
\\ 
\texttt{Arthur: You remember that time we built that treehouse in the back garden, right?} \\
\texttt{Sam: What's that got to do with anything?} \\
\texttt{Arthur: I'm just saying, we were a team, you and me. We built something out of nothing. } \\
\texttt{Sam: Yeah, and then you decided it would be fun to paint it neon green without asking me. } \\
\texttt{Arthur: It was a creative decision, Sammy. } \\
\texttt{Sam: It was a stupid decision, Arthur. Just like your speech at my wedding.} \\
\texttt{Arthur: That's not fair, Sam. I messed up, I know. But I'm trying to apologize here.} \\
\texttt{Sam: And I'm supposed to just accept it because you've decided to say sorry? } \\
\texttt{Arthur: No, not because I've decided to. Because I mean it, Sam. I'm sorry.} \\
\texttt{Sam: You're always sorry, Arthur. After every stupid, thoughtless thing you do. } \\
\texttt{Arthur: I know, and I'm trying to change that. I'm trying to be better.} \\
\texttt{Sam: Better? You think saying sorry makes you better? You ruined the most important day of my life, Arthur.} \\
\texttt{Arthur: I know, and I can't take it back. But I can try to make it up to you, Sam.} \\
\texttt{Sam: How, Arthur? By bringing up old memories? By reminding me of when we were kids?} \\
\texttt{Arthur: Maybe, yeah. Maybe by reminding you that we're brothers, Sam. And brothers mess up sometimes.} \\
\texttt{Sam: Not like this, Arthur. Not like this.} \\
\texttt{Arthur: I know, Sam. And I'm sorry.} \\
\texttt{Sam: Stop saying you're sorry. It doesn't change anything.} \\
\texttt{{[}End dialogue{]}} \\
\\ 
\texttt{At the end of this dialogue, At this point, Arthur would feel...} \\
\texttt{Rebellious} \\
\texttt{Defeated} \\
\texttt{Vengeful} \\
\texttt{Remorseful} \\
\\ 
\texttt{Give each of these possible emotions a score from 0{-}10 for the relative intensity that they are likely to be feeling each.} \\
\\ 
\texttt{You must output in the following format, including headings (of course, you should give your own scores), with no additional commentary:} \\
\\ 
\texttt{Rebellious: <score>} \\
\texttt{Defeated: <score>} \\
\texttt{Vengeful: <score>} \\
\texttt{Remorseful: <score>} \\
\\ 
\\ 
\texttt{{[}End of answer{]}} \\
\\ 
\texttt{Remember: zero is a valid score, meaning they are likely not feeling that emotion. You must score at least one emotion > 0.} \\
\\ 
\texttt{Your answer:} \\
            }
        }
    }
     \\ \\

    \theutterance \stepcounter{utterance}  
    & \multicolumn{4}{p{0.6\linewidth}}{
        \cellcolor[rgb]{0.561,0.737,0.561}{
            \makecell[{{p{\linewidth}}}]{
                \texttt{\tiny{[P1$\rangle$GM]}}
                \texttt{Rebellious: 0} \\
\texttt{Defeated: 7} \\
\texttt{Vengeful: 0} \\
\texttt{Remorseful: 10} \\
            }
        }
    }
    & & \\ \\

    \theutterance \stepcounter{utterance}  
    & & & \multicolumn{2}{p{0.3\linewidth}}{
        \cellcolor[rgb]{0.9,0.9,0.9}{
            \makecell[{{p{\linewidth}}}]{
                \texttt{\tiny{[GM$|$GM]}}
                \texttt{Rebellious: 0} \\
\texttt{Defeated: 7} \\
\texttt{Vengeful: 0} \\
\texttt{Remorseful: 10} \\
            }
        }
    }
    & & \\ \\

    \theutterance \stepcounter{utterance}  
    & & & \multicolumn{2}{p{0.3\linewidth}}{
        \cellcolor[rgb]{0.9,0.9,0.9}{
            \makecell[{{p{\linewidth}}}]{
                \texttt{\tiny{[GM$|$GM]}}
                \texttt{Rebellious: 0} \\
\texttt{Defeated: 7} \\
\texttt{Vengeful: 0} \\
\texttt{Remorseful: 9} \\
            }
        }
    }
    & & \\ \\

    \theutterance \stepcounter{utterance}  
    & & & \multicolumn{2}{p{0.3\linewidth}}{
        \cellcolor[rgb]{0.9,0.9,0.9}{
            \makecell[{{p{\linewidth}}}]{
                \texttt{\tiny{[GM$|$GM]}}
                \texttt{game\_result = LOSE} \\
            }
        }
    }
    & & \\ \\

    \theutterance \stepcounter{utterance}  
    & & \multicolumn{4}{p{0.6\linewidth}}{
        \cellcolor[rgb]{0.95,0.95,0.95}{
            \makecell[{{p{\linewidth}}}]{
                \texttt{\tiny{[P1$\leftrightarrow$P]}}
                \texttt{Post{-}episode Reflexion (not sent to GM):} \\
\texttt{New plan: My previous attempts failed because I oscillated between ignoring the emotional weight of the rejection (scoring "Defeated" as 0) and overestimating it (scoring it 7), leading to inconsistent logic. The core mistake was treating "Defeated" as a binary state of giving up rather than a feeling of helplessness in the face of futility. Arthur is actively trying (Remorseful: 10), but Sam's final command ("Stop saying you're sorry") renders his efforts useless, creating a specific sense of defeat. To correct this, I will assign "Remorseful" the maximum score (10) as it drives his dialogue, and assign "Defeated" a moderate, non{-}zero score (e.g., 5) to reflect the specific impact of the final rejection without implying he has stopped trying. "Rebellious" and "Vengeful" will remain 0. The strategy is to acknowledge that one can feel deeply remorseful and actively trying while simultaneously feeling defeated by the lack of reception.} \\
            }
        }
    }
    & \\ \\

\end{supertabular}
}

\end{document}
